\documentclass[../mathNotesPreamble]{subfiles}

\providecommand{\relscalefact}{1.4}
\begin{document}
\relscale{\relscalefact}
  \section{3.4: Graphs of Polynomial Functions}
    %Characteristics - smooth, continous
    \begin{thmBox*}
      Graphs of polynomials are smooth and continuous with infinite end behavior.
    \end{thmBox*}
    \begin{center}
      \begin{tikzpicture}[
        funcLabel/.style={circle, inner sep=0.5pt, fill=white,opacity=0.0, text opacity=1},
        declare function={
          f(\x)=(\x+2)*\x*(\x-3);
          g(\x)=2*\x-f(3.1);
        }]
        \begin{groupplot}[
          group style={group size=2 by 2,horizontal sep=3cm, vertical sep=3cm},
          axis lines=center,
          axis line style={black,->},
          xmin=-4, xmax=4,
          enlargelimits={abs=0.5},
          xmajorticks=false,
          ymajorticks=false,
          xlabel=, ylabel=,
          ]
          \nextgroupplot[ymin=-8, ymax=8, title={Smooth and continuous}]
            \addplot[<->] expression[domain=-2.5:3.4]{f(x)};
          \nextgroupplot[ymin=-8, ymax=8, title={Not continuous}]
            \addplot[<-] expression[domain=-2.5:3.1]{f(x)};
            \addplot[soldot] coordinates{(3.1,{f(3.1)})};
            \addplot[->] expression[samples at={3.1,3.9}]{g(x)};
            \addplot[holdot] coordinates{(3.1,{g(3.1)})};
          \nextgroupplot[title={Not smooth}, axis equal]
            \addplot[<->] expression[domain=-4:4]{abs(x)};
          \nextgroupplot[title={Finite end behavior}]
            \addplot[<->] expression[domain=-4:4, samples=512]{sin(deg(7*\x))/x};
        \end{groupplot}
      \end{tikzpicture}
    \end{center}
    \pagebreak

    %Find x-intercepts of polynomial functions
    \begin{ex*}
      Find the $x$- and $y$-intercepts of the following functions:
      \begin{extasks}[after-item-skip=\stretch{1}](2)
        \task $f(x)=2(x-4)(x+1)(x-6)$
        \task $g(x)=4x(x-2)^2(x+1)$
        \task $h(x)=x^4-x^2$
        \task $j(x)=x^3+2x^2-9x-18$
      \end{extasks}
      \vspace*{\stretch{1}}
    \end{ex*}
    \pagebreak
    
    %Zeros and multiplicities
    \begin{defn*}[Multiplicity of a Root]
      If a polynomial contains a factor of the form $(x-h)^p$, then the behavior near the $x$-intercept $h$ is determined by the power $p$. We say that $x=h$ is a zero of \textbf{multiplicity} $p$.
      
      When $p$ is, 
      \begin{itemize}
        \item even, the graph will touch the $x$-axis
        \item odd, the graph will cross the $x$-axis
      \end{itemize}
      The degree of the polynomial is the sum of the multiplicities.
    \end{defn*}
    \begin{center}
      \begin{tikzpicture}[
        funcLabel/.style={circle, inner sep=0.5pt, fill=white,opacity=0.0, text opacity=1},
        declare function={
          f(\x)=(\x+10)*\x/10;
          g(\x)=(20+\x)^2*x^2/2500;
          h(\x)=(x+5)*x^3/10;
        }]
        \begin{groupplot}[
          group style={group size=3 by 1,horizontal sep=2.75cm},
          axis lines=center,
          axis line style={black,->},
          xmin=-4, xmax=4,
          ymin=-2, ymax=4,
          enlargelimits={abs=0.5},
          xmajorticks=false,
          ymajorticks=false,
          xlabel=, ylabel=,
          width=0.35\linewidth,
          ]
          \nextgroupplot[title={Single zero}]
            \addplot[<->] expression[domain=-2.5:3.4]{f(x)};
          \nextgroupplot[title={Zero with multiplicity $2$}]
            \addplot[<->] expression[domain=-4:4]{g(x)};
          \nextgroupplot[xmin=-2, xmax=2, title={Zero with multiplicity $3$}]
            \addplot[<->] expression[domain=-1.9:1.8]{h(x)};
        \end{groupplot}
      \end{tikzpicture}
    \end{center}

    \begin{ex*}
      For the following, find the zeros, and the multiplicities of each:
      \begin{extasks}[after-item-skip=\stretch{1}](2)
        \task $f(x)=(x+2)^3(x-3)^2$
        \task $g(x)=x^5(x-1)^3(x+2)$
      \end{extasks}
      \vspace*{\stretch{1}}
    \end{ex*}
    \pagebreak

    %Writing Formulas for Polynomial Functions
    \begin{ex*}
      Find a formula for the polynomial of minimum degree represented in the graph below:
      \begin{center}
        \begin{tikzpicture}[declare function={
          f(\x)=(\x+1)*(\x-1)^3*(\x-4)^2;
          g(\x)=-2*f(\x)/f(0);
          }]
          \begin{axis}[
            axis lines=center,
            axis line style={black,->},
            xtick={-6,-5,...,6},
            xlabel=, ylabel=,
            enlargelimits={abs=0.5},
            width=0.75\linewidth, height=\axisdefaultheight,
            ]
            \addplot[<->] expression[domain=-1.125:4.425]{g(x)};
          \end{axis}
        \end{tikzpicture}
      \end{center}
      \vspace*{\stretch{1}}
    \end{ex*}
    

  \pagebreak
\end{document}
